% =============================================================================
% CompetitiveExam/CS301/11-trees-bst-avl-questions.tex
% Competitive Exam Practice 11 — GATE (Computer Science) Pattern
% Trees: terminology, binary-tree properties, traversals, expression trees,
% BST search/insert/delete, AVL balance factors and rotations, height bounds.
%
% Student-facing — NO answers. Answer key: 11-trees-bst-avl-answers.tex.
%
% Provenance:
%   Q1 [GATE CSE 2009 Q37, verified — ExamSIDE / SolutionsAdda]
%   Q2 [GATE CSE 2020 Q6, verified — ExamSIDE]
%   Q3 [GATE CSE 2017 Set 1, verified — ExamSIDE]
%   Q4 [GATE CSE 2015 Set 1, verified — ExamSIDE]
%   Q5-Q10 [Original, GATE-pattern] — fresh numbers, none reused from the deck,
%     the written assignment, or the coding lab; every answer re-derived.
%
% Compile with: cd CompetitiveExam/CS301 &&
%   TEXINPUTS=../../Preambles:$TEXINPUTS xelatex -interaction=nonstopmode 11-trees-bst-avl-questions.tex
% =============================================================================

\documentclass[11pt]{article}
\usepackage[margin=1in]{geometry}
% =============================================================================
% Preambles/header.tex — shared LaTeX/Beamer preamble for this project
%
% Use from a Beamer lecture by prepending:
%     \input{header}
% and compiling with TEXINPUTS=../Preambles:$TEXINPUTS (the /compile-latex
% skill does this automatically).
%
% PALETTE CONTRACT. Color names defined here MUST match the names in
% Quarto/theme-template.scss so Beamer and Quarto renderings use the same
% palette. The scripts/check-palette-sync.sh script enforces this.
%
% To customize: edit the HEX values below AND the matching $variable in
% Quarto/theme-template.scss, then run scripts/check-palette-sync.sh.
% =============================================================================

% -----------------------------------------------------------------------------
% Engine detection + encoding
%
% Our /compile-latex skill uses XeLaTeX, where inputenc is unnecessary and
% can produce encoding warnings. Only load inputenc under pdfTeX.
% -----------------------------------------------------------------------------
\usepackage{iftex}
\ifPDFTeX
  \usepackage[utf8]{inputenc}
\fi

\usepackage{amsmath, amssymb, amsthm}
\usepackage{booktabs}
\usepackage{graphicx}

% -----------------------------------------------------------------------------
% PALETTE — mirrors Quarto/theme-template.scss variable names.
% Load xcolor and define colors BEFORE anything that references them
% (notably hyperref, hypersetup, and the Beamer color assignments).
% -----------------------------------------------------------------------------
\usepackage{xcolor}

\definecolor{primary-blue}{HTML}{012169}      % headings, accents
\definecolor{primary-gold}{HTML}{B9975B}      % emphasis, borders
\definecolor{highlight-yellow}{HTML}{F2A900}  % markers, alerts
\definecolor{light-bg}{HTML}{E8EDF5}          % title / section backgrounds
\definecolor{jet}{HTML}{1A1A1A}               % body text

% Semantic colors (used in diagrams and annotations)
\definecolor{positive}{HTML}{15803D}          % good / observed / identified
\definecolor{negative}{HTML}{B91C1C}          % bad / problematic / bias
\definecolor{neutral}{HTML}{525252}           % reference / context

% Highlight accents (match .hi-* classes in the SCSS)
\definecolor{hi-slate}{HTML}{314F4F}
\definecolor{hi-green}{HTML}{2E8B57}
\definecolor{hi-red}{HTML}{C41E3A}

% Semantic aliases so skill templates and snippets can write
% \textcolor{accent}{...} without hard-coding "primary-blue".
\colorlet{accent}{primary-blue}
\colorlet{accent2}{primary-gold}

% -----------------------------------------------------------------------------
% Hyperref — loaded AFTER xcolor + palette so link colors resolve.
% -----------------------------------------------------------------------------
\usepackage{hyperref}
\hypersetup{colorlinks=true, linkcolor=primary-blue, urlcolor=primary-blue,
            citecolor=primary-blue, pdfborder={0 0 0}}

% -----------------------------------------------------------------------------
% Course code — set once per deck (after \input{header}, before \begin{document})
% via \coursecode{CS401}. Shown in the footer on every slide so a deck's course
% is identifiable without opening the title page. Empty by default.
% -----------------------------------------------------------------------------
\makeatletter
\newcommand{\coursecode}[1]{\gdef\@coursecodetext{#1}}
\gdef\@coursecodetext{}
\makeatother

% -----------------------------------------------------------------------------
% Beamer theme assignments (only apply under Beamer).
%
% We detect Beamer via \@ifclassloaded{beamer}, which is the documented,
% non-internal way. This must be wrapped in \makeatletter/\makeatother
% because \@ifclassloaded contains an @-catcode character.
% -----------------------------------------------------------------------------
\makeatletter
\@ifclassloaded{beamer}{%
  \usefonttheme{professionalfonts}
  \setbeamercolor{structure}{fg=primary-blue}
  \setbeamercolor{frametitle}{fg=primary-blue}
  \setbeamercolor{title}{fg=primary-blue}
  \setbeamercolor{subtitle}{fg=primary-gold}
  \setbeamercolor{author}{fg=primary-blue}
  \setbeamercolor{institute}{fg=neutral}
  \setbeamercolor{date}{fg=primary-gold}
  \setbeamercolor{itemize item}{fg=primary-blue}
  \setbeamercolor{itemize subitem}{fg=primary-gold}
  \setbeamercolor{alerted text}{fg=primary-gold}
  \setbeamercolor{block title}{fg=white, bg=primary-blue}
  \setbeamercolor{block body}{bg=light-bg}
  % Semantic block variants: block=definition/concept (blue), exampleblock=
  % worked example (green/positive), alertblock=key takeaway or caution (gold).
  % Keep to <=2 colored boxes per slide across ALL three variants (INV-7).
  \setbeamercolor{block title example}{fg=white, bg=positive}
  \setbeamercolor{block body example}{bg=light-bg}
  \setbeamercolor{block title alerted}{fg=white, bg=primary-gold}
  \setbeamercolor{block body alerted}{bg=light-bg}

  % Minimal footer: course code (left) + page X / N (right), no navigation clutter
  \setbeamertemplate{navigation symbols}{}
  \setbeamertemplate{footline}{%
    \color{neutral}\footnotesize\hspace{0.7em}\@coursecodetext\hfill
    \insertframenumber/\inserttotalframenumber\hspace{0.7em}\vspace{0.3em}}
}{}
\makeatother

% -----------------------------------------------------------------------------
% TikZ setup — libraries the snippet gallery uses
% -----------------------------------------------------------------------------
\usepackage{tikz}
\usetikzlibrary{arrows.meta, positioning, calc, decorations.pathreplacing,
                fit, shapes.geometric, backgrounds}

% Shared TikZ styles — snippets and hand-written diagrams can pick these up
% via `\begin{tikzpicture}[every node/.style={...}]` or by naming specific
% styles (e.g., `\node[dag-node] (X) at (0,0) {X};`).
\tikzset{
  % Node styles for causal DAGs and similar
  dag-node/.style = {
    draw=primary-blue, fill=light-bg, rounded corners=2pt,
    minimum width=1.2cm, minimum height=0.9cm, align=center, font=\small
  },
  decision-node/.style = {
    draw=primary-gold, fill=white, diamond, aspect=1.8,
    minimum width=1.8cm, minimum height=1.2cm, align=center, font=\small,
    inner sep=1pt
  },
  % Edge styles — observed vs counterfactual
  observed-edge/.style = {-{Latex[length=2.5mm]}, thick, draw=jet},
  counterfactual-edge/.style = {-{Latex[length=2.5mm]}, thick, draw=neutral, dashed},
  confound-edge/.style = {-{Latex[length=2.5mm]}, thick, draw=negative, dashed},
  % Data-point styles for plots
  observed-dot/.style = {fill=primary-blue, draw=primary-blue, circle, inner sep=1.2pt},
  counterfactual-dot/.style = {fill=white, draw=neutral, circle, inner sep=1.2pt}
}

% -----------------------------------------------------------------------------
% Convenience macros
% -----------------------------------------------------------------------------
\newcommand{\muted}[1]{\textcolor{neutral}{#1}}
\newcommand{\key}[1]{\textcolor{primary-gold}{\textbf{#1}}}
\newcommand{\good}[1]{\textcolor{positive}{#1}}
\newcommand{\bad}[1]{\textcolor{negative}{#1}}

% Standout frame macro: full-bleed dark background for section transitions.
% Beamer-only (uses \begin{frame}/\setbeamercolor/\usebeamerfont) -- do not
% call from an article-class file (e.g. Notes/<CODE>/*-notes.tex). \muted,
% \key, \good, \bad, and \coursecode above are plain xcolor/gdef and are
% safe to use from either class; verified when Notes/article-class support
% was added.
% Expand: \transitionslide{Your transition title}
\newcommand{\transitionslide}[1]{%
  \begingroup
  \setbeamercolor{background canvas}{bg=primary-blue}
  \begin{frame}[plain]
    \centering
    \vfill
    {\usebeamerfont{title}\color{white}\Large #1\par}
    \vfill
  \end{frame}
  \endgroup
}

% Section divider frame (Beamer-only), an alternative to \transitionslide with a
% darker two-line style: near-black (jet) background, a small white label line
% above a larger gold title, and a thin gold rule along the bottom edge.
% Expand: \sectiondivider{Section 2}{Pointers}
\newcommand{\sectiondivider}[2]{%
  \begingroup
  \setbeamercolor{background canvas}{bg=jet}
  \begin{frame}[plain]
    \vspace*{3.0cm}
    {\color{white}\ttfamily\large #1\par}
    \vspace{0.5cm}
    {\color{primary-gold}\LARGE #2\par}
    \begin{tikzpicture}[remember picture, overlay]
      \draw[primary-gold, line width=2.5pt]
        ($(current page.south west)+(0,0.1)$) -- ($(current page.south east)+(0,0.1)$);
    \end{tikzpicture}
  \end{frame}
  \endgroup
}

% -----------------------------------------------------------------------------
% Channel watermark — "Concepts That Click" (YouTube).
%
% Article-class files (CompetitiveExam Q/A sets, Notes, Assignments) get a
% light diagonal draft watermark on every page plus a centered footer naming
% the channel. Beamer decks get a subtle TikZ background watermark instead
% (the footline already carries the course code). Centralized here so every
% MCQ PDF is branded without per-file edits.
% -----------------------------------------------------------------------------
\makeatletter
\@ifclassloaded{article}{%
  \usepackage{draftwatermark}
  \SetWatermarkText{Concepts That Click}
  \SetWatermarkScale{1}
  \SetWatermarkFontSize{2cm}
  \SetWatermarkLightness{0.86}
  \SetWatermarkAngle{45}
  \usepackage{fancyhdr}
  \pagestyle{fancy}
  \fancyhf{}
  \fancyfoot[C]{\footnotesize\textcolor{neutral}{Concepts That Click~|~YouTube: \href{https://www.youtube.com/@concepts.thatclick}{@concepts.thatclick}}}
  \renewcommand{\headrulewidth}{0pt}
  \renewcommand{\footrulewidth}{0pt}
  \fancypagestyle{plain}{%
    \fancyhf{}%
    \fancyfoot[C]{\footnotesize\textcolor{neutral}{Concepts That Click~|~YouTube: \href{https://www.youtube.com/@concepts.thatclick}{@concepts.thatclick}}}%
    \renewcommand{\headrulewidth}{0pt}%
    \renewcommand{\footrulewidth}{0pt}%
  }
}{}
\@ifclassloaded{beamer}{%
  \setbeamertemplate{background}{%
    \begin{tikzpicture}[remember picture, overlay]
      \node[rotate=45, opacity=0.055, align=center]
        at (current page.center)
        {\Huge\textcolor{neutral}{Concepts That Click}};
    \end{tikzpicture}%
  }
}{}
\makeatother 
\coursecode{CS301}
\renewcommand{\thesection}{11.\arabic{section}}
\setlength{\emergencystretch}{3em}

\title{Competitive Exam Practice 11: Trees, BSTs and AVL\\\large GATE (Computer Science) Pattern}
\author{Auqib Hamid Lone\\[0.3em]
  \emph{Department of Computer Science \& Engineering}, \emph{GCET Ganderbal}}
\date{\today}

\begin{document}
\maketitle

\noindent Each question carries 1--2 marks. MCQ = single correct option; NAT = Numerical Answer Type. Height is measured in \emph{edges} unless a question states otherwise: a single node has height $0$, the empty tree $-1$. A BST satisfies, at every node, left subtree $<$ node $<$ right subtree; the balance factor is $\mathrm{height}(\text{left})-\mathrm{height}(\text{right})$.

\begin{enumerate}
\item[\textbf{Q1}] \textit{[GATE CSE 2009 Q37, verified --- ExamSIDE / SolutionsAdda]} \textbf{[MCQ, 1 mark]} What is the maximum height of any AVL-tree with $7$ nodes? Assume that the height of a tree with a single node is $0$.
\begin{enumerate}
  \item[(A)] $2$
  \item[(B)] $3$
  \item[(C)] $4$
  \item[(D)] $5$
\end{enumerate}

\item[\textbf{Q2}] \textit{[GATE CSE 2020 Q6, verified --- ExamSIDE]} \textbf{[MCQ, 1 mark]} What is the worst case time complexity of inserting $n^{2}$ elements into an AVL-tree with $n$ elements initially?
\begin{enumerate}
  \item[(A)] $\Theta(n^{4})$
  \item[(B)] $\Theta(n^{2})$
  \item[(C)] $\Theta(n^{2}\log n)$
  \item[(D)] $\Theta(n^{3})$
\end{enumerate}

\item[\textbf{Q3}] \textit{[GATE CSE 2017 Set 1, verified --- ExamSIDE]} \textbf{[MCQ, 1 mark]} Let $T$ be a binary search tree with $15$ nodes. The minimum and maximum possible heights of $T$ are:
\emph{Note:} The height of a tree with a single node is $0$.
\begin{enumerate}
  \item[(A)] $4$ and $15$ respectively
  \item[(B)] $3$ and $14$ respectively
  \item[(C)] $4$ and $14$ respectively
  \item[(D)] $3$ and $15$ respectively
\end{enumerate}

\item[\textbf{Q4}] \textit{[GATE CSE 2015 Set 1, verified --- ExamSIDE]} \textbf{[MCQ, 1 mark]} The height of a tree is the length of the longest root-to-leaf path in it. The maximum and minimum number of nodes in a binary tree of height $5$ are
\begin{enumerate}
  \item[(A)] $63$ and $6$, respectively
  \item[(B)] $64$ and $5$, respectively
  \item[(C)] $32$ and $6$, respectively
  \item[(D)] $31$ and $5$, respectively
\end{enumerate}

\item[\textbf{Q5}] \textit{[Original, GATE-pattern]} \textbf{[MCQ, 2 marks]} A binary tree has inorder traversal
\[
  D,\ B,\ E,\ A,\ F,\ C,\ G
\]
and preorder traversal
\[
  A,\ B,\ D,\ E,\ C,\ F,\ G.
\]
Which one of the following is its postorder traversal?
\begin{enumerate}
  \item[(A)] $D,\ E,\ B,\ F,\ G,\ C,\ A$
  \item[(B)] $D,\ B,\ E,\ A,\ F,\ C,\ G$
  \item[(C)] $A,\ B,\ D,\ E,\ C,\ F,\ G$
  \item[(D)] $D,\ E,\ B,\ C,\ F,\ G,\ A$
\end{enumerate}

\item[\textbf{Q6}] \textit{[Original, GATE-pattern]} \textbf{[MCQ, 2 marks]} Keys
\[
  50,\ 30,\ 70,\ 20,\ 40,\ 60,\ 80
\]
are inserted in that order into an empty BST. The key $30$ is then deleted using the inorder-successor rule. Which one of the following is the resulting inorder traversal?
\begin{enumerate}
  \item[(A)] $20,\ 40,\ 50,\ 60,\ 70,\ 80$
  \item[(B)] $20,\ 30,\ 40,\ 50,\ 60,\ 70,\ 80$
  \item[(C)] $20,\ 40,\ 50,\ 60,\ 80,\ 70$
  \item[(D)] $40,\ 20,\ 50,\ 60,\ 70,\ 80$
\end{enumerate}

\item[\textbf{Q7}] \textit{[Original, GATE-pattern]} \textbf{[MCQ, 2 marks]} Keys $40,10,20$ are inserted in that order into an initially empty AVL tree. How many rotations are performed, and what is the final root?
\begin{enumerate}
  \item[(A)] one rotation; root $10$
  \item[(B)] one rotation; root $20$
  \item[(C)] two rotations; root $20$
  \item[(D)] two rotations; root $40$
\end{enumerate}

\item[\textbf{Q8}] \textit{[Original, GATE-pattern]} \textbf{[NAT, 2 marks]} The tree below is a BST.
\[
50;\quad 30\ (20,\ 40\ (\,\_\,,35));\quad 70\ (60).
\]
Give the balance factor of node $30$, the balance factor of node $40$, and the height of the whole tree, in that order. Is the tree an AVL tree? \underline{\hspace{4cm}}

\item[\textbf{Q9}] \textit{[Original, GATE-pattern]} \textbf{[NAT, 1 mark]} Let $N(h)$ be the minimum number of nodes in an AVL tree of height $h$ (a single node has height $0$). What is $N(4)$? \underline{\hspace{2cm}}

\item[\textbf{Q10}] \textit{[Original, GATE-pattern]} \textbf{[MCQ, 2 marks]} The postfix expression
\[
  4\ 6\ +\ 2\ *\ 7\ -
\]
is represented by an expression tree. What is the value of the expression, and which traversal of the tree prints the original postfix form?
\begin{enumerate}
  \item[(A)] $13$; postorder
  \item[(B)] $13$; preorder
  \item[(C)] $21$; postorder
  \item[(D)] $21$; inorder
\end{enumerate}
\end{enumerate}

\end{document}
